% -----------------------------------------------------------------------------------------------------
\section{单圈尖点反常维度}\label{sec:appB-cusp}

\begin{proposition}[单圈尖点反常维度 \tagR]\label{prop:appB-cusp-oneloop}
（$R_{\mathrm{PT}}$：在重整化微扰论中对红外正规化的部分子态严格。）在 $O(\alpha_s)$，对散射组态的折线
$W_\beta=\frac1{N_c}\langle0|\Tr\,T\{W[\infty v_1,0]\,W[0,-\infty v_2]\}|0\rangle$（$v_i^2=1$，$v_i^0>0$，$\cosh\beta=v_1\cdot v_2$；Feynman 规范，$D=4-2\varepsilon$，红外由绝热阻尼 $\ee^{-\delta|\lambda|}$ 正规化），
\begin{equation}\label{eq:appB-cusp-pole}
W_\beta=1-\frac{\alpha_sC_F}{2\pi\varepsilon}\,(\beta\coth\beta-1)+(\varepsilon\text{ 的正则项})+O(\alpha_s^2),\qquad
\Gamma_{\rm cusp}(\beta)=\frac{\alpha_sC_F}{\pi}\,(\beta\coth\beta-1)+O(\alpha_s^2),
% VF-CONV-23
\end{equation}
其中 $\Gamma_{\rm cusp}$ 按\cref{def:sudakov-cusp}定义（$\MSbar$）。$\beta\coth\beta$ 来自两线之间的交换，$-1$ 来自两条线的自能；$\Gamma(0)=0$；
$\beta\to\infty$ 时 $\Gamma=\frac{\alpha_sC_F}\pi(\beta-1)+O(\beta\ee^{-2\beta})$，线性项的系数 $\frac{\alpha_sC_F}\pi=\frac{\alpha_s}{4\pi}\Gamma_0$ 即单圈的类光 cusp 反常维度，
而 $v_i=(\cosh y_i,0,0,\sinh y_i)$ 时 $\beta=|y_1-y_2|$。
\uses{prop:appB-rules,def:prelim-dr-conventions,def:prelim-msbar,def:prelim-feynman-rules,def:sudakov-cusp}
\end{proposition}

\begin{proof}
\emph{第一步：图。}$O(g^2)$ 有三个图：两条线之间的单胶子交换，与每条线自身的 $O(g^2)$ 项（自能）。由\cref{prop:appB-rules}的单胶子项与 Schwinger 表示
$\frac{-\ii}{v\cdot k-\ii\delta}=\int_0^\infty\dd l\,\ee^{-\ii l(v\cdot k-\ii\delta)}$、$\frac{\ii}{v\cdot k+\ii\delta}=\int_{-\infty}^0\dd\lambda\,\ee^{-\ii\lambda(v\cdot k+\ii\delta)}$，
$W^{(1)}[\infty v_1,0]=\ii g\int_0^\infty\dd l\,\ee^{-\delta l}\,v_1\cdot A(lv_1)$，$W^{(1)}[0,-\infty v_2]=\ii g\int_{-\infty}^0\dd\lambda\,\ee^{\delta\lambda}\,v_2\cdot A(\lambda v_2)$。
Feynman 规范中 $\langle0|TA^a_\mu(x)A^b_\nu(y)|0\rangle=-g_{\mu\nu}\delta^{ab}D_F(x-y)$（\cref{def:prelim-feynman-rules}），
$D_F(x)=\int\frac{\dd^Dk}{(2\pi)^D}\frac{\ii\,\ee^{-\ii k\cdot x}}{k^2+\ii0}=\frac{\Gamma(1-\varepsilon)}{4\pi^{2-\varepsilon}}(-x^2+\ii0)^{-1+\varepsilon}$（Wick 转动后的欧氏 Fourier 变换，\cref{def:prelim-dr-conventions}）。
色因子 $\frac1{N_c}\Tr(t^at^a)=C_F$。令 $s=-\lambda$，
\[
\begin{aligned}
E&=g^2C_F\cosh\beta\int_0^\infty\!\dd l\int_0^\infty\!\dd s\,\ee^{-\delta(l+s)}D_F(lv_1+sv_2),\\
\Sigma_i&=g^2C_F\!\!\int\limits_{0<l_1<l_2}\!\!\dd l_1\dd l_2\,\ee^{-\delta(l_1+l_2)}D_F\big((l_2-l_1)v_i\big)=\frac{g^2C_F}{2\delta}\int_0^\infty\!\dd\tau\,\ee^{-\delta\tau}D_F(\tau v_i)
\end{aligned}
\]
（$E$ 的符号：$(\ii g)^2\cdot(-v_1\cdot v_2)=g^2\cosh\beta$；$\Sigma_i$ 中 $(\ii g)^2\cdot(-v_i^2)=g^2$，令 $\tau=l_2-l_1$ 后对 $l_1$ 积分得 $1/(2\delta)$）。

\emph{第二步：交换图。}令 $l=\rho u$、$s=\rho(1-u)$（$\dd l\,\dd s=\rho\,\dd\rho\,\dd u$），$(lv_1+sv_2)^2=\rho^2N(u)$，$N(u)=u^2+(1-u)^2+2u(1-u)\cosh\beta=1+2u(1-u)(\cosh\beta-1)\ge1$。
间隔是类时的，$(-\rho^2N+\ii0)^{-1+\varepsilon}=-\ee^{\ii\pi\varepsilon}(\rho^2N)^{-1+\varepsilon}$，于是
\[
\begin{aligned}
E&=-g^2C_F\cosh\beta\,\frac{\ee^{\ii\pi\varepsilon}\Gamma(1-\varepsilon)}{4\pi^{2-\varepsilon}}\int_0^\infty\!\dd\rho\,\rho^{-1+2\varepsilon}\ee^{-\delta\rho}\int_0^1\!\dd u\,N(u)^{-1+\varepsilon}\\
&=-g^2C_F\cosh\beta\,\frac{\ee^{\ii\pi\varepsilon}\Gamma(1-\varepsilon)\Gamma(2\varepsilon)}{4\pi^{2-\varepsilon}\,\delta^{2\varepsilon}}\int_0^1\!\dd u\,N(u)^{-1+\varepsilon}.
\end{aligned}
\]
$\rho\to0$ 处的发散（尖点处的紫外发散）给出极点 $\Gamma(2\varepsilon)=\frac1{2\varepsilon}+O(1)$。令 $w=2u-1$，$N=\frac12(a^2-b^2w^2)$，$a^2=\cosh\beta+1$，$b^2=\cosh\beta-1$，$ab=\sinh\beta$，
\[
\int_0^1\frac{\dd u}{N(u)}=\int_{-1}^{1}\frac{\dd w}{a^2-b^2w^2}=\frac{1}{ab}\ln\frac{a+b}{a-b}=\frac{\beta}{\sinh\beta}\qquad\Big(\frac{a+b}{a-b}=\frac{\cosh\frac\beta2+\sinh\frac\beta2}{\cosh\frac\beta2-\sinh\frac\beta2}=\ee^\beta\Big).
\]
故 $E$ 的极点为 $-\frac{g^2C_F}{8\pi^2\varepsilon}\beta\coth\beta=-\frac{\alpha_sC_F}{2\pi\varepsilon}\beta\coth\beta$。

\emph{第三步：自能。}$D_F(\tau v)=-\ee^{\ii\pi\varepsilon}\frac{\Gamma(1-\varepsilon)}{4\pi^{2-\varepsilon}}\tau^{-2+2\varepsilon}$，而
$\int_0^\infty\dd\tau\,\tau^{-2+2\varepsilon}\ee^{-\delta\tau}=\Gamma(-1+2\varepsilon)\,\delta^{1-2\varepsilon}$
（在 $\Real\varepsilon>\frac12$ 收敛，再解析延拓：线性紫外发散在维数正规化中不留痕迹），故
\[
\Sigma_i=-g^2C_F\,\frac{\ee^{\ii\pi\varepsilon}\Gamma(1-\varepsilon)\Gamma(-1+2\varepsilon)}{8\pi^{2-\varepsilon}\,\delta^{2\varepsilon}} .
\]
由 $\Gamma(-1+2\varepsilon)=-\Gamma(2\varepsilon)/(1-2\varepsilon)=-\frac1{2\varepsilon}+O(1)$，
每条线的自能极点为 $+\frac{g^2C_F}{16\pi^2\varepsilon}=+\frac{\alpha_sC_F}{4\pi\varepsilon}$；它是半无限线在有限端点（尖点）处的对数紫外发散。

\emph{第四步：反常维度。}代入 $g_0^2=4\pi\alpha_s\mu^{2\varepsilon}\ee^{\varepsilon\gamma_E}(4\pi)^{-\varepsilon}$（\cref{def:prelim-msbar}），$\mu^{2\varepsilon}\delta^{-2\varepsilon}=(\mu^2/\delta^2)^\varepsilon$，
三图之和的极点为 $-\frac{\alpha_sC_F}{2\pi\varepsilon}(\beta\coth\beta-1)$，即式~\eqref{eq:appB-cusp-pole} 的第一式；故 $Z_{\rm cusp}=1-\frac{\alpha_sC_F}{2\pi\varepsilon}(\beta\coth\beta-1)$，
由 $\frac{\dd\alpha_s}{\dd\ln\mu}=-2\varepsilon\alpha_s+O(\alpha_s^2)$，$\Gamma_{\rm cusp}=\frac{\dd\ln Z_{\rm cusp}}{\dd\ln\mu}=\frac{\alpha_sC_F}\pi(\beta\coth\beta-1)$。
等价地，$W^{\rm ren}_\beta=1-\frac{\alpha_sC_F}{2\pi}(\beta\coth\beta-1)\ln\frac{\mu^2}{\delta^2}+\text{常数}$，$-\frac{\dd\ln W^{\rm ren}_\beta}{\dd\ln\mu}=\Gamma_{\rm cusp}$。

\emph{第五步：极限。}$\beta\coth\beta-1=\frac{\beta^2}3+O(\beta^4)$，$\Gamma(0)=0$：直线没有尖点发散，交换图的 $\int\dd u/N=1$ 恰好抵消两段在连接点的端点发散；
$\beta\to\infty$ 时 $\coth\beta=1+2\ee^{-2\beta}+\cdots$；$v_1\cdot v_2=\cosh(y_1-y_2)$。各积分与常数因子都经符号与数值计算核对。% check: s5_cusp.py
\end{proof}

\paragraph{两圈系数与文献注记。}在 $\Gamma_{\rm cusp}=\sum_n(\frac{\alpha_s}{4\pi})^{n+1}\Gamma_n$ 的归一化下，类光两圈系数为
$\Gamma_1=4C_F[(\frac{67}9-\frac{\pi^2}3)C_A-\frac{20}9T_Fn_f]$\cite{KorchemskyRadyushkin1987}，$n_f=5$ 时 $\Gamma_1=36.8436$；\cite[式~(A.10)]{BecherSchwartz2008}（JHEP 版与 arXiv v1、v2）在 $T_Fn_f$ 前多印了一个 $C_F$，按字面得到 $26.97$（附录~Q 第~1~条）。% VF-CONV-23, VF-SCET-16
式~\eqref{eq:appB-cusp-pole} 包含两条线的自能；Becher 与 Neubert 的有质量部分子形式\cite[式~(13)--(15)]{Becher:2009kw}只把 $\beta\coth\beta$ 记入尖点项、把 $-1$ 归入 $2\gamma^Q$（$\gamma^Q_0=-2C_F$）。
单圈结果见\cite{Polyakov1980,BrandtNeriSato1981}；产生组态（$\beta\to\beta-\ii\pi$）见第~5~章\cref{prop:sudakov-coulomb}。
